
(@@M0@@ \uline{et} @@M1@@
\uline{étant affectés du poids} @@M2@@), \uline{et ces} @@M3@@ \uline{relations engendrent l'idéal des relations entre les} @@M4@@.
\end{enumerate}

\textbf{Corollaire 4.4.} \uline{Les éléments de} @@M5@@ \uline{de la forme}~:
@@M6@@
\uline{forment une base de} @@M7@@ \uline{considéré comme module sur} @@M8@@.

Pour la démonstration du lemme et du corollaire, nous renvoyons à [2], exposé 4, pages 4.19, auquel ils sont empruntés.

4.5.
Soient @@M9@@ un schéma quasi-compact, @@M10@@ un @@M11@@-Module localement libre de rang @@M12@@, @@M13@@ une suite d'entiers @@M14@@, de somme @@M15@@. Posons @@M16@@ ; @@M17@@ est muni d'une suite de quotients canonique @@M18@@ ; on pose @@M19@@, @@M20@@, @@M21@@ étant le morphisme structural @@M22@@. Pour @@M23@@ on note @@M24@@ le noyau de @@M25@@. Soit @@M26@@ le groupe de Grothendieck des faisceaux localement libres de type fini sur @@M27@@ (resp. des complexes parfaits). On note @@M28@@ la classe de @@M29@@, et @@M30@@ la classe de @@M31@@. Dans le cas où @@M32@@ est défini en termes de complexes parfaits, on donne encore un sens évident aux séries @@M33@@, @@M34@@, où @@M35@@ est la classe de @@M36@@. Enfin, on remarque que~:
@@M37@@

\textbf{Théorème 4.6.} \uline{Avec les hypothèses et les notations de} 4.5, \uline{la} @@M38@@-\uline{algèbre} @@M39@@ \uline{est engendrée par les} @@M40@@; \uline{l'idéal des relations entre les} @@M41@@ \uline{est engendré par les relations déduites de}~:
@@M42@@
